\documentclass[a4paper,fleqn]{article}
\usepackage{amsfonts}
\usepackage{amsmath}
\usepackage{amssymb}
\usepackage{graphicx}     

\setlength{\mathindent}{0pt} % Zorgt ervoor dat wiskundige expressies helemaal links staan

\begin{document}
  
\noindent \large Auteur: Victoria Mazur \\
\noindent \large Studierichting: Informatica\\
\noindent \large Oefeningenreeks 1D nr. 19.12\\

\medskip

\normalsize

\begin{enumerate}
    \item Oefening:
    $$
    z^5 + 2z^4 + 2z^3 + 4z^2 - 15z - 30
    $$

    \item Ontbinding in factoren:
    $$
    z^4(z + 2) + 2z^2(z + 2) - 15(z + 2) = 0
    $$
    $$
    (z^4 + 2z^2 - 15)(z + 2) = 0
    $$
    \begin{itemize}
        \item \( z + 2 = 0 \Rightarrow z = -2 \)
        \item \( z^4 + 2z^2 - 15 = 0 \)
    \end{itemize}

    \item Substitueer \( w = z^2 \) in \( z^4 + 2z^2 - 15 = 0 \):
    $$
    w^2 + 2w - 15 = 0
    $$
    $$
    (w + 5)(w - 3) = 0
    $$
    Dit geeft:
    \begin{itemize}
        \item \( w + 5 = 0 \Rightarrow w = -5 \Rightarrow z^2 = -5 \Rightarrow z = \pm i\sqrt{5} \)
        \item \( w - 3 = 0 \Rightarrow w = 3 \Rightarrow z^2 = 3 \Rightarrow z = \pm \sqrt{3} \)
    \end{itemize}

    \item Nulpunten in \(\mathbb{C}\):
    $$
    z = -2, \quad z = i\sqrt{5}, \quad z = -i\sqrt{5}, \quad z = \sqrt{3}, \quad z = -\sqrt{3}
    $$

    \item Eindresultaat \(\mathbb{R}[z]\):
    $$
    (z + 2)(z^2 + 5)(z + \sqrt{3})(z - \sqrt{3})
    $$
\end{enumerate}

\begin{thebibliography}{999}
%\bibitem{a} Referentie
\end{thebibliography}
\end{document}


